\documentclass[11pt]{amsart}

\usepackage[T1]{fontenc}
\usepackage{lmodern}
\usepackage{microtype}
\usepackage{amsmath,amssymb,amsthm}
\usepackage[hidelinks]{hyperref}

\hypersetup{
  pdftitle={A Characteristic-Two Counterexample to the Shalev--Smoktunowicz Conjecture},
  pdfsubject={Finite braces and powerful 2-groups},
  pdfkeywords={left brace, powerful 2-group, right nilpotence, regular subgroup}
}

\newtheorem{theorem}{Theorem}
\newtheorem{lemma}[theorem]{Lemma}
\newtheorem{proposition}[theorem]{Proposition}
\newtheorem{corollary}[theorem]{Corollary}

\newcommand{\F}{\mathbb{F}}
\newcommand{\Span}{\operatorname{span}}
\newcommand{\Soc}{\operatorname{Soc}}

\title[A characteristic-two counterexample]
{A Characteristic-Two Counterexample to the\\
Shalev--Smoktunowicz Conjecture}

\date{}

\subjclass[2020]{16T25, 20D15, 16N40}
\keywords{left brace, powerful $2$-group, right nilpotence, regular subgroup}

\begin{document}

\begin{abstract}
We construct, for every $r\geq1$, a left brace $A_r$ of order $2^{8r+1}$
whose additive group is elementary abelian and whose multiplicative group is
a powerful $2$-group of nilpotency class two, while $A_r$ is not right
nilpotent. The case $r=1$ has order $512$ and multiplicative group $G$
satisfying $G'=G^4\cong C_2^2$ and $\operatorname{exp}(G)=8$. This gives a
characteristic-two counterexample to the Shalev--Smoktunowicz conjecture and
answers Question~10.2 of Verbeken. Together with Verbeken's odd-prime
examples, it completes the negative answer for every prime. The construction
is a characteristic-two adaptation of Verbeken's kernel-graph method; the
essential modification is a truncation degree divisible by four.
\end{abstract}

\maketitle

\section{Introduction}

For a left brace $A$, write
\[
  \lambda_a(b)=-a+a\circ b,
  \qquad
  a*b=\lambda_a(b)-b.
\]
For additive subgroups $X,Y\leq(A,+)$, let $X*Y$ denote the additive
subgroup generated by the elements $x*y$. The left and right series are
\[
 A^1=A,\quad A^{n+1}=A*A^n,
 \qquad
 A^{(1)}=A,\quad A^{(n+1)}=A^{(n)}*A.
\]
A finite $2$-group $G$ is powerful when
$G'\leq G^4=\langle g^4:g\in G\rangle$ \cite{LubotzkyMann}.

Shalev and Smoktunowicz asked whether a brace of prime-power order with
powerful adjoint group must be strongly nilpotent
\cite[Question~1]{ShalevSmoktunowicz}. Vendramin recorded the corresponding
right-nilpotence statement as Problem~5.15
\cite[Problem~5.15]{Vendramin}. Following Verbeken, we refer to this
right-nilpotence statement as the Shalev--Smoktunowicz conjecture.
Verbeken disproved it for every odd prime and asked whether an analogous
counterexample exists for $p=2$ \cite[Question~10.2]{Verbeken}.
Together with his odd-prime examples, Theorem~\ref{thm:main} settles the
conjecture negatively for every prime.

Our construction is the characteristic-two instance of Verbeken's
kernel-graph mechanism \cite[Proposition~2.2]{Verbeken}. The new point is to
decouple the truncation degree from the characteristic and take it divisible
by four; this places the commutators inside fourth powers, as required for a
powerful $2$-group.

\begin{theorem}\label{thm:main}
For every $r\geq1$ there exists a left brace $A_r$ such that
\begin{enumerate}
  \item $(A_r,+)\cong C_2^{\,8r+1}$;
  \item if $G_r=(A_r,\circ)$, then $G_r$ has nilpotency class two,
        $G_r'\cong C_2^2$, and $G_r'\leq G_r^4$;
  \item $A_r$ is left nilpotent, is not right nilpotent, and
        $\Soc(A_r)=0$.
\end{enumerate}
For $r=1$, one has
\[
 |A_1|=512,
 \qquad
 G_1'=G_1^4\cong C_2^2,
 \qquad
 \operatorname{exp}(G_1)=8.
\]
\end{theorem}

\begin{corollary}\label{cor:order}
The order-$3^7$ example in \cite{Verbeken} is not minimal among finite left
braces with powerful multiplicative group that are not right nilpotent.
\end{corollary}

\begin{proof}
Theorem~\ref{thm:main} gives an example of order
$2^9=512<3^7=2187$.
\end{proof}

No claim is made that $512$ is the smallest possible order.

\section{The construction}

Fix $r\geq1$, put $m=4r$, and work over $\F_2$. Let
\[
 B=\F_2[x,y]/(xy,x^{m+1}+y^{m+1}),
 \qquad
 J=(x,y),
 \qquad
 z=x^{m+1}=y^{m+1}.
\]

\begin{lemma}\label{lem:algebra}
The ideal $J$ has basis
\[
 x,x^2,\ldots,x^m,
 y,y^2,\ldots,y^m,z.
\]
In particular, $\dim_{\F_2}J=2m+1$ and
$J^{m+1}=\F_2z$, $J^{m+2}=0$.
The assignments
\[
 D(x)=y^m,
 \qquad
 D(y)=x^m
\]
define a derivation $D$ of $B$. If
\[
 W=D(J)=\Span_{\F_2}\{x^m,y^m\},
\]
then $D^2=0$, $W^2=0$, and $\phi=1+D$ is an algebra automorphism of order
two.
\end{lemma}

\begin{proof}
With lexicographic order $x>y$, the polynomials
\[
 xy,\qquad x^{m+1}+y^{m+1},\qquad y^{m+2}
\]
form a Gr\"obner basis of the defining ideal. The standard monomials are
$1,x,\ldots,x^m,y,\ldots,y^{m+1}$, which proves the basis assertion and the
statements about the powers of $J$.

The derivation preserves the relations because
\[
 D(xy)=y^{m+1}+x^{m+1}=0,
 \qquad
 D(x^{m+1}+y^{m+1})=0.
\]
Moreover, $D(x^i)=D(y^i)=0$ for every $i\geq2$. Hence $D^2=0$ and
$D(J)=W$, while $W^2=0$. Therefore
\[
 (a+D(a))(b+D(b))=ab+D(ab)
\]
for all $a,b\in B$, so $\phi=1+D$ is multiplicative. Since $\phi^2=1$ and
$D\neq0$, it has order two.
\end{proof}

Let $U=1+J\leq B^\times$.

\begin{lemma}\label{lem:character}
There is a homomorphism $\ell:U\to\F_2$ such that
\[
 \ell(1+x)=\ell(1+y)=0,
 \qquad
 \ell(1+z)=1,
 \qquad
 \ell\circ\phi=\ell.
\]
\end{lemma}

\begin{proof}
Since $U$ is abelian, $U/U^2$ is an $\F_2$-vector space. The classes of
$1+x$, $1+y$, and $1+z$ are linearly independent. Indeed,
\[
 (1+x)^\alpha(1+y)^\beta(1+z)^\gamma
 =1+\alpha x+\beta y+\gamma z,
\]
whereas every element of $U^2$ is a square $1+a^2$ and has no $x$-, $y$-,
or $z$-component; the last assertion uses that $m+1$ is odd. A linear
functional on $U/U^2$ with the prescribed values therefore lifts to a
homomorphism $\ell:U\to\F_2$.

For $a\in J$, write $a=\alpha x+\beta y+a_2$, where $a_2$ has no linear
terms. Then $D(a_2)=0$, $a_2W=0$, and
\[
 D(a)=\alpha y^m+\beta x^m,
 \qquad
 aD(a)=\alpha\beta(x^{m+1}+y^{m+1})=0.
\]
Thus, for $u=1+a$,
\[
 \phi(u)=u(1+D(a)).
\]
Because $m=4r$,
\[
 1+D(a)=(1+\alpha y^r+\beta x^r)^4\in U^4\subseteq U^2.
\]
Every homomorphism from $U$ to $\F_2$ vanishes on $U^2$, so
$\ell(\phi(u))=\ell(u)$.
\end{proof}

Put $\varepsilon(a)=\ell(1+a)$ and define, for $a,b\in J$,
\begin{equation}\label{eq:circle}
 \lambda_a(b)=(1+a)\phi^{\varepsilon(a)}(b),
 \qquad
 a\circ b=a+\lambda_a(b).
\end{equation}

\begin{proposition}\label{prop:brace}
The operations $+$ and $\circ$ make $J$ a left brace, denoted $A_r$. Its
star product is
\begin{equation}\label{eq:star}
 a*b=ab+\varepsilon(a)D(b)+\varepsilon(a)aD(b).
\end{equation}
\end{proposition}

\begin{proof}
This is the kernel-graph construction of
\cite[Proposition~2.2]{Verbeken}; we include the short verification. The map
\[
 \chi:U\rtimes\langle\phi\rangle\longrightarrow\F_2,
 \qquad
 \chi(u\phi^t)=\ell(u)+t,
\]
is a homomorphism by Lemma~\ref{lem:character}. Its kernel is
\[
 H=\{u\phi^{\ell(u)}:u\in U\}.
\]
Acting on $U$ by
\[
 (u\phi^t)\mathbin{\cdot}v=u\,\phi^t(v),
\]
the group $H$ is regular because the element indexed by $u$ sends $1$ to
$u$. Under the identification $1+a\leftrightarrow a$, its affine action on
$J$ is
\[
 b\longmapsto a+(1+a)\phi^{\varepsilon(a)}(b).
\]
The standard regular-subgroup construction therefore gives the left-brace
operation \eqref{eq:circle}. Since
$\phi^{\varepsilon(a)}=1+\varepsilon(a)D$, expansion gives
\eqref{eq:star}.
\end{proof}

\section{Proof of the theorem}

Let $A=A_r$ and $G=(A,\circ)$, identified with the kernel graph above.
Lemma~\ref{lem:algebra} gives
$(A,+)\cong C_2^{\,2m+1}=C_2^{\,8r+1}$. Put
$K=\ker\ell$. Since $\ell(1+z)=1$, the element
$c=(1+z)\phi$ belongs to $G$, and
\[
 G=K\langle c\rangle.
\]
Conjugation by $c$ induces $\phi$ on the abelian group $K$. Hence
\[
 G'=\langle k^{-1}\phi(k):k\in K\rangle.
\]
For $k=1+a$, Lemma~\ref{lem:character} gives
$k^{-1}\phi(k)=1+D(a)$, so $G'\leq1+W$. Since
$1+x,1+y\in K$,
\[
 (1+x)^{-1}\phi(1+x)=1+y^m,
 \qquad
 (1+y)^{-1}\phi(1+y)=1+x^m.
\]
Therefore
\begin{equation}\label{eq:derived}
 G'=1+W\cong C_2^2.
\end{equation}
The subgroup $1+W$ is fixed pointwise by $\phi$ and is central in $U$,
hence central in $G$. Thus $G$ has class at most two; by
\eqref{eq:derived}, its class is exactly two.

Let $g=u\phi^t\in G$, where $u=1+a$ and $t=\ell(u)$. Then
\[
 g^2=u^2(1+D(a))^t,
 \qquad
 g^4=u^4.
\]
Consequently,
\begin{equation}\label{eq:fourth}
 \begin{split}
 G^4=U^4
 &=1+\operatorname{im}(a\mapsto a^4)\\
 &=1+\Span_{\F_2}\{x^4,x^8,\ldots,x^m,
                         y^4,y^8,\ldots,y^m\}.
 \end{split}
\end{equation}
In particular, $G'=1+W\leq G^4$, so $G$ is powerful. If $r=1$, then
\eqref{eq:fourth} gives $G^4=1+W=G'$. Moreover, every fourth power then
lies in the exponent-two group $1+W$, so $\operatorname{exp}(G)\leq8$;
the element $1+x\in G$ (which lies in $K$) has order eight, since
$(1+x)^4=1+x^4\neq1$ and $(1+x)^8=1$. Hence
$\operatorname{exp}(G)=8$.

We next examine the right series. Let
\[
 T=\Span_{\F_2}\{x^m,y^m,z\}.
\]
Since $1+x^m,1+y^m\in U^4$, one has
$\varepsilon(x^m)=\varepsilon(y^m)=0$, while $\varepsilon(z)=1$. Formula
\eqref{eq:star} gives
\[
 x^m*x=z,
 \qquad
 y^m*y=z,
 \qquad
 z*x=y^m,
 \qquad
 z*y=x^m.
\]
Also $TJ\subseteq\F_2z$, $D(J)=W$, and $TW=0$, so
$T*A\subseteq T$. The displayed products show the reverse inclusion and
also show $T\subseteq A^{(2)}$. Thus
\[
 T*A=T,
 \qquad
 T\subseteq A^{(n)}\quad(n\geq2),
\]
and $A$ is not right nilpotent.

For the left series, \eqref{eq:star} gives $A*A\subseteq J^2$. Since
$\varepsilon(x)=\varepsilon(y)=0$,
\[
 x*x^i=x^{i+1},
 \qquad
 y*y^i=y^{i+1}
 \qquad(1\leq i\leq m),
\]
so $A*A=J^2$. Moreover, $D(J^n)=0$ for $n\geq2$, and therefore
$A*J^n=J^{n+1}$. Inductively,
\[
 A^n=J^n,
 \qquad
 A^{m+1}=\F_2z,
 \qquad
 A^{m+2}=0.
\]
Hence $A$ is left nilpotent.

Finally, suppose $a\in\Soc(A)=\ker\lambda$ and put
$e=\varepsilon(a)$. From \eqref{eq:circle},
\[
 \lambda_a(x)=x+e y^m+ax+eay^m.
\]
Neither $ax$ nor $ay^m$ has a $y^m$-component, so
$\lambda_a(x)=x$ implies $e=0$. It follows that
$(1+a)b=b$ for all $b\in J$, hence $aJ=0$. The basis in
Lemma~\ref{lem:algebra} gives
$\operatorname{Ann}_J(J)=\F_2z$, so $a=\gamma z$. But then
$0=e=\varepsilon(a)=\gamma$, and therefore $a=0$. This proves
$\Soc(A)=0$ and completes the proof of Theorem~\ref{thm:main}.

\begin{thebibliography}{9}

\bibitem{LubotzkyMann}
A.~Lubotzky and A.~Mann,
\emph{Powerful $p$-groups. I. Finite groups},
J. Algebra \textbf{105} (1987), no.~2, 484--505.
\href{https://doi.org/10.1016/0021-8693(87)90211-0}
{doi:10.1016/0021-8693(87)90211-0}.

\bibitem{ShalevSmoktunowicz}
A.~Shalev and A.~Smoktunowicz,
\emph{Groups, Lie rings and braces},
Bull. Belg. Math. Soc. Simon Stevin \textbf{31} (2024), no.~4, 422--433.
\href{https://doi.org/10.36045/j.bbms.230929}
{doi:10.36045/j.bbms.230929}.

\bibitem{Vendramin}
L.~Vendramin,
\emph{Skew braces: a brief survey},
in \emph{Geometric Methods in Physics XL}, Trends in Mathematics,
Birkh\"auser/Springer, Cham, 2024, 153--175.
\href{https://doi.org/10.1007/978-3-031-62407-0_12}
{doi:10.1007/978-3-031-62407-0\_12}.

\bibitem{Verbeken}
B.~Verbeken,
\emph{Powerful multiplicative groups do not force right nilpotence in finite braces},
arXiv:2608.01884v1 [math.GR], 2026.
\href{https://arxiv.org/abs/2608.01884}{arXiv:2608.01884}.

\end{thebibliography}

\end{document}